\section*{Chapitre \ref{ch:mx:portfolio-trades-and-risk-bids} --- Transactions de portefeuille et risk bids}

\begin{solution}{exo:mx:portfolio-trades-and-risk-bids:1}
$0.3/0.2=1.5$ jours~; $3+0.7\times200\times\sqrt{0.2}=65.6$ points de base.
\end{solution}

\begin{solution}{exo:mx:portfolio-trades-and-risk-bids:2}
La variance vaut $\sigma^2d/3$~: $150\sqrt{3/3}=150$ points de base.
\end{solution}

\begin{solution}{exo:mx:portfolio-trades-and-risk-bids:3}
$4\times\E[\max\text{ de 6 normales}]=4\times1.27=5.07$ points de base.
\end{solution}

\begin{solution}{exo:mx:portfolio-trades-and-risk-bids:4}
Un panier long de 500 actions relève surtout du risque de marché, que l'on supprime d'un coup par une seule transaction sur contrats à terme. Il reste le risque sectoriel et de style (les autres
facteurs) ainsi que le risque spécifique des actions, que le contrat à terme ne couvre pas~; avec 500 titres, la part spécifique est faible, de sorte que les autres facteurs dominent.
\end{solution}

\begin{solution}{exo:mx:portfolio-trades-and-risk-bids:5}
Le fonds détenait, dans chaque tranche, davantage de titres liquides qu'un tirage aléatoire dans la tranche n'en donnerait~; le profil ne permet pas de le voir. Un panier penché vers les titres
moins liquides de chaque tranche serait sous-estimé, ce qui constitue le danger réel de l'enchérisseur à l'aveugle.
\end{solution}

\begin{solution}{exo:mx:portfolio-trades-and-risk-bids:6}
Les erreurs deviennent corrélées~: la part commune de l'erreur n'est pas sélectionnée par l'enchère (tous les enchérisseurs la partagent), de sorte que la malédiction issue de la part indépendante diminue,
mais tous les enchérisseurs se trompent ensemble, et le client gagne ou perd en fonction de l'erreur commune.
\end{solution}

\begin{solution}{exo:mx:portfolio-trades-and-risk-bids:7}
Coût 14.9, charge de risque 25.0, décote 6.2 (l'erreur d'estimation croît avec la base plus grande), marge 2.0~: 48.1 points de base.
\end{solution}

\begin{solution}{exo:mx:portfolio-trades-and-risk-bids:8}
Gagner signifie que les cours étaient bas par rapport aux concurrents, non qu'ils couvraient les coûts~: la malédiction du vainqueur sélectionne les transactions dont les coûts
ont été sous-estimés. Seul le P\&L réalisé des transactions remportées, comparé à leurs cours, dit si le prix était juste.
\end{solution}

\begin{solution}{pb:mx:portfolio-trades-and-risk-bids:1}
\textbf{1.} Agence~: le client supporte les coûts et le risque~; principal~: le courtier prend les positions à un prix convenu.
\textbf{2.} Voir les définitions.
\textbf{3.} 786 titres, volume quotidien médian de 37.5 millions de dollars, volatilité médiane de 171 points de base~; un panier de 250 millions de dollars de 500 titres, position médiane
de 1.0~\% du volume quotidien.
\textbf{4.} Une économie sur les coûts de transaction pour le client et une source efficace de flux pour les fournisseurs de liquidité.
\textbf{5.} Au plus 20~\% du volume quotidien, demi-fourchette plus $Y\sigma\sqrt{\min(Q/V,0.2)}$~; 14.9 points de base.
\textbf{6.} $\int_0^a(1-t/a)(1-t/b)\,dt=a-a/2-a^2/(2b)+a^2/(3b)=a/2-a^2/(6b)$.
\textbf{7.} 40.4 points de base sans couverture, 12.5 avec couverture.
\textbf{8.} 98.8~\% du notionnel est inférieur à 5~\% du volume quotidien~; aucune position n'exige plus d'une journée.
\textbf{9.} La plus basse de quatre estimations indépendantes est trop basse de $1.03s$ en moyenne~; il faut décoter de ce montant.
\textbf{10.} $-2.2$ points de base sans décote, 2.0 (la marge) avec décote.
\textbf{11.} En tarifant 300 paniers au même profil~: coût 17.0 (écart type 0.4), risque couvert 13.0 (1.1).
\textbf{12.} Panier annoncé~: 14.9, 12.5, 4.2 et 2.0~; à l'aveugle~: 17.0, 13.0, 4.4 et 2.0.
\textbf{13.} \emph{Résultat nommé}~: 33.6 points de base sous la clôture pour le panier annoncé et 36.4 pour le panier à l'aveugle~; la décote
liée à la malédiction du vainqueur face à trois concurrents vaut 4.2 et 4.4 points de base.
\textbf{14.} 50.4, 32.8 et 24.0 points de base du fonds.
\textbf{15.} Elle ne paie aucune fourchette et ne déplace aucun cours~: les deux parties sont satisfaites au point milieu.
\textbf{16.} Lorsque les noms risqueraient de fuir vers des dealers qui opéreraient par anticipation sur le débouclage d'un montant supérieur à la prime.
\textbf{17.} Le coefficient du modèle d'impact et les erreurs des concurrents sont supposés indépendants.
\textbf{18.} Décoter moins pour l'erreur commune, mais élargir la marge pour le risque de modèle.
\textbf{19.} Le principal achète la certitude à un prix~; l'agence est moins chère en moyenne lorsque le client peut supporter le risque.
\textbf{20.} La certitude de la clôture, en contrepartie du coût et du risque du débouclage et de la malédiction du vainqueur.
\end{solution}

\begin{solution}{iq:mx:portfolio-trades-and-risk-bids:1}
Les noms (ou le profil), les tailles rapportées au volume, les fourchettes et les volatilités, la covariance et les couvertures disponibles, l'indice de référence, la concurrence. Prix
= coût de liquidation attendu + une charge pour le risque du débouclage + décote liée à la malédiction du vainqueur + marge.
\iqlookfor{Les quatre composantes~; couverture.}
\end{solution}

\begin{solution}{iq:mx:portfolio-trades-and-risk-bids:2}
L'estimation gagnante est la plus optimiste de plusieurs estimations bruitées~; il faut corriger en ajoutant l'avance attendue de la meilleure de $n$ estimations, $s\E[\max]$,
ou en tarifant conditionnellement au fait de gagner.
\iqlookfor{Sélection sur l'estimation la plus basse.}
\end{solution}

\begin{solution}{iq:mx:portfolio-trades-and-risk-bids:3}
La variance du P\&L pendant le débouclage, $\sum w_iw_jC_{ij}\int x_ix_j\,dt$ avec un modèle de risque~; il faut couvrir le marché (et les secteurs ou les styles) par des contrats à terme
ou des ETF d'emblée, ce qui laisse le risque spécifique et le risque factoriel non couvrable.
\iqlookfor{Covariance pondérée par le temps~; couvertures.}
\end{solution}

\begin{solution}{iq:mx:portfolio-trades-and-risk-bids:4}
Savoir si le profil masque des inclinaisons défavorables (les paniers du client sont moins liquides dans les tranches qu'on ne le suppose), si la décote reflète le nombre réel
de concurrents et l'erreur du modèle, et si les pertes proviennent du coût, du risque ou de la fuite après l'attribution.
\iqlookfor{Biais du profil~; décote~; décomposition du P\&L.}
\end{solution}

\begin{solution}{iq:mx:portfolio-trades-and-risk-bids:5}
Inventorier les deux portefeuilles, compenser les titres communs, croiser en interne, couvrir l'exposition hors marché par des contrats à terme, planifier le reste selon la liquidité, et
rapporter le coût à l'estimation préalable à la transaction.
\iqlookfor{Compensation, croisement, couverture, reporting.}
\end{solution}

\begin{solution}{iq:mx:portfolio-trades-and-risk-bids:6}
Tirer des paniers conformes au profil et les tarifer~; vérifier le biais en tarifant après coup, comme s'ils étaient à l'aveugle, les paniers annoncés et
les comparer à leur coût réel sur de nombreuses transactions.
\iqlookfor{Simulation à partir du profil~; test rétrospectif de l'estimation à l'aveugle.}
\end{solution}
